%! Absolute Value Functions \& Transformations — Unit 2, Lesson 2.3 — Guided Notes
\documentclass[10pt]{article}
\usepackage{algebra2-article}
\usepackage{algebra2-boxes}
\newcommand{\TallMath}[1]{$\displaystyle #1\rule[-1.4em]{0pt}{3.2em}$}
% Coordinate grid: {xmin}{xmax}{ymin}{ymax}
\newcommand{\lgrid}[4]{%
  \draw[step=1, linegray!40, very thin] (#1,#3) grid (#2,#4);
  \draw[->, semithick, charcoal] (#1,0)--(#2,0) node[below right, font=\scriptsize]{$x$};
  \draw[->, semithick, charcoal] (0,#3)--(0,#4) node[left, font=\scriptsize]{$y$};}

\begin{document}

\pageheader{Unit 2, Lesson 2.3}{Guided Notes}

\vspace{0.10in}

\begin{objectivebox}
\textbf{By the end of this lesson, I will be able to\dots}
\begin{itemize}\setlength{\itemsep}{2pt}
  \item graph the \emph{absolute value parent} $f(x)=|x|$ and read its vertex, axis, and minimum;
  \item use \emph{vertex form} $g(x)=a|x-h|+k$ to find the vertex $(h,k)$ and axis $x=h$;
  \item decide from $a$ whether a V \emph{opens up or down} (min vs.\ max) and is narrower or wider;
  \item describe any absolute value graph as \emph{shifts, a stretch, and a reflection} of $y=|x|$.
\end{itemize}
\end{objectivebox}

\vspace{0.05in}

\begin{vocabbox}
\small
Fill in each term as we build it together.
\par\vspace{2pt}   % \par is required: \termblanklong's \noindent is a no-op mid-paragraph
\termblanklong{Absolute value parent function $f(x)=|x|$}
\termblanklong{Vertex (turning point)}
\termblanklong{Axis of symmetry}
\termblanklong{Vertex form $a|x-h|+k$}
\termblanklong{Reflection / stretch / compression}
\end{vocabbox}

\begin{hookbox}
You drive toward a rest stop, reach it after $2$ hours, and keep going. Your \emph{distance from the
rest stop} (in tens of miles) is $d(t)=30\,|t-2|$.
\begin{itemize}\setlength{\itemsep}{1pt}
  \item Distance at $t=0$? \blank{1.4cm} \quad At $t=2$? \blank{1.4cm} \quad At $t=4$? \blank{1.4cm}
  \item You are \emph{closest} at $t=\blank{1.0cm}$, where the distance is \blank{1.0cm}.
\end{itemize}
Plot those points and the shape is a \textbf{V}. A distance is never negative, so the graph bottoms
out at one turning point --- the \emph{vertex} --- and is symmetric about the line through it. Every
absolute value function is this same V, moved and stretched.
\end{hookbox}

\boxguard[18]
\begin{notesbox}{1. The Parent Function $f(x)=|x|$}
Complete the table, then read the graph.
\begin{center}
\renewcommand{\arraystretch}{1.15}
\begin{minipage}{0.40\linewidth}
\centering
\begin{tabular}{|c|c|}
\hline
$x$ & $y=|x|$ \\ \hline
$-2$ & \blank{0.9cm} \\ \hline
$-1$ & \blank{0.9cm} \\ \hline
$0$ & \blank{0.9cm} \\ \hline
$1$ & \blank{0.9cm} \\ \hline
$2$ & \blank{0.9cm} \\ \hline
\end{tabular}
\end{minipage}%
\begin{minipage}{0.55\linewidth}
\centering
\begin{tikzpicture}[scale=0.5]
  \lgrid{-4}{4}{-1}{5}
  \draw[very thick, forest] (-4,4)--(0,0)--(4,4);
  \fill[charcoal] (0,0) circle (3.5pt);
\end{tikzpicture}
\end{minipage}
\end{center}
\textbf{Read the V:} vertex $(\blank{0.7cm},\,\blank{0.7cm})$; \; axis of symmetry
$x=\blank{0.8cm}$; \; minimum value \blank{0.8cm}; \; domain \blank{2.0cm}; \; range \blank{1.6cm}.
It \emph{decreases} on \blank{1.8cm} and \emph{increases} on \blank{1.8cm}. The V is really two lines
pasted at the vertex: $y=\blank{0.8cm}$ when $x\ge 0$, and $y=\blank{0.8cm}$ when $x<0$.
\end{notesbox}

\boxguard[26]
\begin{notesbox}{2. Up / Down and Left / Right}
\textbf{Outside} the bars acts \emph{vertically, as expected}; \textbf{inside} the bars acts
\emph{horizontally and backwards}.
\begin{center}
\renewcommand{\arraystretch}{1.3}
\begin{tabularx}{\linewidth}{@{}l l X@{}}
$y=|x|+k$ & shifts the V \blank{2.2cm} & $y=|x|-3$: vertex $(\blank{0.6cm},\,\blank{0.6cm})$ \\
$y=|x-h|$ & shifts the V \blank{2.2cm} & $y=|x-3|$: vertex $(\blank{0.6cm},\,\blank{0.6cm})$; \quad $y=|x+2|$: vertex $(\blank{0.6cm},\,\blank{0.6cm})$ \\
\end{tabularx}
\end{center}
Each solid V below is the parent (dashed) after \emph{one} shift. Name each vertex on the line.
\begin{center}
\begin{tikzpicture}[scale=0.44]
  % left: y = |x - 3|  (right 3)
  \begin{scope}
    \lgrid{-2}{6}{-1}{5}
    \draw[very thick, forest!40, dashed] (-1,1)--(0,0)--(4,4);
    \draw[very thick, forest] (-1,4)--(3,0)--(6,3);
    \fill[charcoal] (3,0) circle (3.5pt);
    \node[charcoal, font=\scriptsize] at (2,-2.4){$y=|x-3|$};
  \end{scope}
  % right: y = |x| - 3  (down 3)
  \begin{scope}[xshift=11cm]
    \lgrid{-4}{4}{-4}{3}
    \draw[very thick, forest!40, dashed] (-3,3)--(0,0)--(3,3);
    \draw[very thick, forest] (-3,0)--(0,-3)--(3,0);
    \fill[charcoal] (0,-3) circle (3.5pt);
    \node[charcoal, font=\scriptsize] at (0,-5.2){$y=|x|-3$};
  \end{scope}
\end{tikzpicture}
\end{center}
Left vertex: $(\blank{0.7cm},\,\blank{0.7cm})$. \qquad Right vertex: $(\blank{0.7cm},\,\blank{0.7cm})$.
\end{notesbox}

\boxguard[20]
\begin{notesbox}{3. Stretch, Compress, and Reflect: $y=a|x|$}
The number $a$ multiplies the \emph{outputs}: it scales the V and, if negative, flips it.
\begin{center}
\begin{tikzpicture}[scale=0.42]
  % y = 2|x|  (narrow)
  \begin{scope}
    \lgrid{-3}{3}{-1}{5}
    \draw[very thick, forest!40, dashed] (-3,3)--(0,0)--(3,3);
    \draw[very thick, forest] (-2.2,4.4)--(0,0)--(2.2,4.4);
    \node[charcoal, font=\scriptsize] at (0,-2.4){$y=2|x|$};
  \end{scope}
  % y = 1/2 |x|  (wide)
  \begin{scope}[xshift=9cm]
    \lgrid{-4}{4}{-1}{5}
    \draw[very thick, forest!40, dashed] (-3,3)--(0,0)--(3,3);
    \draw[very thick, forest] (-4,2)--(0,0)--(4,2);
    \node[charcoal, font=\scriptsize] at (0,-2.4){$y=\tfrac12|x|$};
  \end{scope}
  % y = -|x|  (down)
  \begin{scope}[xshift=19cm]
    \lgrid{-3}{3}{-4}{2}
    \draw[very thick, forest!40, dashed] (-2,2)--(0,0)--(2,2);
    \draw[very thick, forest] (-3,-3)--(0,0)--(3,-3);
    \node[charcoal, font=\scriptsize] at (0,-5.4){$y=-|x|$};
  \end{scope}
\end{tikzpicture}
\end{center}
\begin{itemize}[leftmargin=1.5em, itemsep=1pt]
  \item $|a|>1$ makes the V \blank{2.4cm}; \quad $0<|a|<1$ makes it \blank{2.4cm}.
  \item $a<0$ \blank{2.6cm} the V, so it opens \blank{1.6cm} instead of up.
  \item The value of $a$ is the \emph{slope} of the \blank{2.6cm} ray of the V.
\end{itemize}
\end{notesbox}

\boxguard[23]
\begin{notesbox}{4. Vertex Form: $g(x)=a|x-h|+k$}
All three moves at once. The vertex is $(h,k)$ and the axis of symmetry is $x=h$.
\[ \text{opens \textbf{up} if } a>0 \ (k \text{ is a } \textbf{minimum}); \qquad
   \text{opens \textbf{down} if } a<0 \ (k \text{ is a } \textbf{maximum}). \]
\textbf{Worked:} $g(x)=-2|x-1|+3$.
\begin{center}
\renewcommand{\arraystretch}{1.3}
\begin{tabularx}{\linewidth}{@{}X X X X@{}}
Vertex $(\blank{0.6cm},\,\blank{0.6cm})$ & Axis $x=\blank{0.8cm}$ & Opens \blank{1.4cm} & \blank{1.2cm} value $=\blank{0.8cm}$ \\
\end{tabularx}
\end{center}
Because $a=-2$, the V is \blank{1.6cm} (narrower / wider) than the parent and opens
\blank{1.4cm}. Evaluate at $x=0$ straight from the rule:
\begin{work}
  g(0) &= -2\,|\,0-1\,| + 3 \\
       &= -2\,|{-1}| + 3     \\
       &= -2(1) + 3          \\
       &= 1
\end{work}
\end{notesbox}

\begin{practicebox}
For $g(x) = \tfrac12\,|x+4| - 2$, fill in each feature.
\begin{enumerate}[leftmargin=1.7em, itemsep=3pt]
  \item Vertex $(\blank{0.8cm},\,\blank{0.8cm})$; \quad axis of symmetry $x=\blank{1.0cm}$.
  \item Opens \blank{1.4cm} (up / down), so $k$ is a \blank{1.2cm} (max / min) value of \blank{0.8cm}.
  \item Compared to $y=|x|$, this graph is shifted \blank{1.6cm} $4$ and \blank{1.6cm} $2$, and is
        \blank{1.4cm} (narrower / wider).
  \item Evaluate $g(0)$, one step per line.
        \begin{work}
          g(0) &= \tfrac12\,|0+4| - 2 \\
               &= \tfrac12(4) - 2      \\
               &= 2 - 2                \\
               &= 0
        \end{work}
\end{enumerate}
\end{practicebox}

\end{document}
